
\doublespacing

\title{\Large\textbf{Earthquakes and Divergent Economic Recovery:\\
Case Studies of Chile and New Zealand}}

\author[1]{Diego A. Díaz}
\author[2]{Pablo Paniagua}
\author[1]{Cristián Larroulet}

\affil[1]{School of Business and Economics (FEN), Universidad del Desarrollo}
\affil[2]{Department of Political Economy, King's College London}


\maketitle

\begin{abstract}
\singlespacing
The consequences of natural disasters, such as earthquakes, are evident: death, coordination problems, destruction of infrastructure, and displacement of the population. However, according to empirical research, the impact of a natural disaster on economic activity is mixed. This is relevant for highly seismic countries such as Chile and New Zealand. This paper contributes to the literature on natural disasters and their economic effects by analyzing the cases of two affected regions within these countries. We employ the synthetic control method (SCM) to estimate the medium-term impact of two large earthquakes on regional output: the 2010 Maule (Chile) event and the 2010--2011 Canterbury (New~Zealand) sequence. We find that Chile and New Zealand experienced opposite economic effects: the GDP per capita of the Canterbury region rose above its synthetic counterfactual by about 12\% in the years following the disaster, while the GDP per capita of the Maule region showed no significant change relative to its counterfactual. In interpreting these outcomes, we argue that New~Zealand’s insurance architecture, implementation capacity, and community co-production enabled a rapid, large-scale rebuild, while Chile’s institutional context supported recovery back to trend but not above trend. The paper highlights that \textit{institutional readiness} mediates disaster impacts and underscores that the measured gains are medium-lived output responses.
\end{abstract}

